  All non-canonically isomorphic to
  $(k+h^\vee) H^3(\fg)[[k+h^\vee]]$.
\end{enumerate}
\end{theorem}

\begin{theorem}[Deformation families of $\Ethree$-algebras]
\label{thm:e3-identification}
Let $\fg$ be a simple finite-dimensional Lie algebra and
$V_k(\fg)$ the affine Kac--Moody vertex algebra at
level~$k$.
Writing $\lambda = k + h^\vee$ for the departure from
critical level, the derived chiral centre
$Z^{\mathrm{der}}_{\mathrm{ch}}(V_k(\fg))$ of
Proposition~\textup{\ref{prop:e3-structure}} and the
CFG $\Ethree$-algebra $\cA^\lambda$ of
Theorem~\textup{\ref{thm:cfg}} are isomorphic as formal
deformation families of $\Ethree$-algebras:
\begin{equation}\label{eq:e3-families}
  \{Z^{\mathrm{der}}_{\mathrm{ch}}(V_k(\fg))
  \}_{k + h^\vee \in \lambda\CC[[\lambda]]}
  \;\simeq\; \{\cA^\lambda\}_{\lambda
  \in \lambda H^3(\fg)[[\lambda]]}
\end{equation}
over the base $\lambda H^3(\fg)[[\lambda]]$
\textup{(}with $\lambda = k + h^\vee$\textup{)}.
\begin{enumerate}[label=\textup{(\roman*)}]
\item \textup{(Matching deformation spaces.)}
  Both families are parametrised by
  $(k+h^\vee) H^3(\fg)[[k+h^\vee]]$:
  for the derived chiral centre, the deformation
  parameter is the departure from critical level;
  for the CFG algebra, it is the Chern--Simons coupling.
\item \textup{(Base point mismatch.)}
  At the critical level $k = -h^\vee$, the derived
  chiral centre $Z^{\mathrm{der}}_{\mathrm{ch}}
  (V_{-h^\vee}(\fg))$ is related to the
  Feigin--Frenkel centre
  $\mathrm{Fun}(\mathrm{Op}_{\fg^\vee}(D))$
  \textup{(}opers on the formal disk\textup{)}, not
  directly to $C^*(\fg)$. The isomorphism is
  perturbative: it holds in the formal neighbourhood
  of critical level, order by order in
  $\lambda = k + h^\vee$.
\item \textup{(Order-by-order uniqueness.)}
  For simple $\fg$, $H^3(\fg) \cong \CC$ is
  one-dimensional, so the deformation base is
  one-dimensional at each order in $\lambda$. The
  isomorphism~\eqref{eq:e3-families} is forced by
  the matching of associated graded algebras
  \textup{(}both $C^*(\fg)$\textup{)},
  matching $\Pthree$ brackets on the formal disk
  \textup{(}Theorem~\textup{\ref{thm:chiral-e3-cfg}(ii)}%
  \textup{)}, and the one-dimensionality of the
  obstruction space at each order.
\item \textup{(Sugawara constraint.)}
  The topological enhancement
  \textup{(}Theorem~\ref{thm:e3-cs}(iv)\textup{)} holds
  at generic $k$ but fails at $k = -h^\vee$, where
  the Sugawara element is undefined. The perturbative
  identification holds
  order by order in $(k + h^\vee)$, since the Sugawara
  element exists at any nonzero departure from critical
  level.
\end{enumerate}
\end{theorem}

\begin{proof}
The argument is obstruction-theoretic: we construct an
isomorphism of $\Ethree$-algebras order by order in
$\lambda = k + h^\vee$, using the one-dimensionality of
$H^3(\fg)$ for simple~$\fg$ to force uniqueness at each
stage.

\smallskip
\noindent\textit{Step 1: setup.}
Both $Z^{\mathrm{der}}_{\mathrm{ch}}(V_k(\fg))$ and
$\cA^\lambda$ are filtered $\Ethree$-algebras deforming
$C^*(\fg) = \Sym(\fg^\vee[-1])$ over the base
$\lambda H^3(\fg)[[\lambda]]$, with $\lambda = k + h^\vee$
(Theorem~\ref{thm:e3-cs}(ii) and
Theorem~\ref{thm:cfg}(iv) respectively).
Both filtrations are complete and exhaustive with
associated graded $C^*(\fg)$.
The deformation theory of $\Ethree$-algebras with fixed
associated graded $C^*(\fg)$ is controlled by the
$\Ethree$-deformation complex, whose relevant cohomology
group is $H^3(\fg)$: the space of $\Ethree$-deformations
of $C^*(\fg)$ at each order $p \geq 1$ in $\lambda$,
modulo those at order $< p$, is classified by
$H^3(\fg)$ (Lurie~\cite{HA},
Costello--Francis--Gwilliam~\cite{CFG26},
Theorem~1.4).

\smallskip
\noindent\textit{Step 2: induction on order.}
We construct the isomorphism
$\varphi \colon Z^{\mathrm{der}}_{\mathrm{ch}}(V_k(\fg))
\xrightarrow{\sim} \cA^\lambda$ as a compatible family of
isomorphisms $\varphi_p$ modulo $\lambda^{p+1}$, by
induction on~$p$.

\emph{Base case ($p = 0$):} both associated graded
algebras are $C^*(\fg)$
(Theorem~\ref{thm:chiral-e3-cfg}(i)), so
$\varphi_0 = \mathrm{id}_{C^*(\fg)}$.

\emph{Inductive step:} suppose $\varphi_p$ is an
$\Ethree$-isomorphism modulo $\lambda^{p+1}$.
The obstruction to extending $\varphi_p$ to an
isomorphism modulo $\lambda^{p+2}$ lies in the
space of $\Ethree$-deformations of $C^*(\fg)$ at
order $p+1$, which is $H^3(\fg)$. For simple~$\fg$,
$H^3(\fg) = \CC$ is one-dimensional, generated by the
Cartan three-form
$\omega_3(X, Y, Z) = (X, [Y, Z])$.
Two $\Ethree$-deformations of $C^*(\fg)$
that agree modulo $\lambda^{p+1}$ therefore differ at order
$p+1$ by a scalar multiple of $\omega_3$.

It remains to show that both constructions select the
\emph{same} scalar at each order. On the formal disk
$D = \Spec \CC[[z]]$, the restriction functor
$\Gamma(D, -)$ is symmetric monoidal and preserves
$\Ethree$-structures
(Theorem~\ref{thm:chiral-e3-cfg}, Step~1 of the
proof). By Theorem~\ref{thm:chiral-e3-cfg}(ii), the
chiral $\Pthree$ bracket restricted to~$D$ recovers
the CFG $\Pthree$ bracket at every order. Since the
$\Pthree$ bracket determines the $\Ethree$-deformation
class (the bracket on generators
$\fg^*[-1] \otimes \fg^*[-1] \to \CE^{\mathrm{ch}}(\fg_k)$
encodes the full deformation datum for biderivations, by
the argument of
Lemma~\ref{lem:bv-p3-commutativity}, Step~2), and the
space of $\fg$-equivariant corrections at each order is
one-dimensional (spanned by the Killing-form term
$k\,(a,b)\,\partial$), the two constructions agree at
order $p+1$. The isomorphism $\varphi_{p+1}$ extends
$\varphi_p$.

\smallskip
\noindent\textit{Step 3: passage to the limit.}
The compatible family $\{\varphi_p\}_{p \geq 0}$
determines an isomorphism
$\varphi = \varprojlim_p\,\varphi_p$
of filtered $\Ethree$-algebras over
$\lambda\,H^3(\fg)[[\lambda]]$, completing the proof
of~\eqref{eq:e3-families}.
Part~(iv) is immediate: the topological enhancement
(Theorem~\ref{thm:e3-cs}(iv)) requires the Sugawara
element, which exists for $\lambda \neq 0$ and is
compatible with the isomorphism~$\varphi$ order by order
(the Sugawara element at order~$p$ depends only on the
data modulo~$\lambda^{p+1}$, where the two families
already agree).
\end{proof}

% ----------------------------------------------------------------
\subsection{Explicit $\Ethree$ operations for
$V_k(\mathfrak{sl}_2)$}
\label{subsec:e3-explicit}
